% Input time-series profiles (load factor, energy price) over the T = 6 horizon.
% Styling follows input_profiles_fig.tex of the tADMM paper exactly: same
% colours, same axis geometry, same legend placement, same left/right axis
% split. There is no PV module in this benchmark, so that series is absent.
%
% B&W readability uses THREE independent channels per curve so the figure
% survives a grayscale printout: (i) colour (Rose Pine), (ii) line style,
% (iii) marker shape -- shape is luminance-independent and disambiguates even
% when tones collapse to similar greys.
%   Load factor  : pine       / solid    / filled circle
%   Energy price : cash green / dash-dot / filled triangle
%
% Unlike the tADMM paper, the curves are NOT analytic shapes sampled at T
% points. The data are given at the nodes; the curves are the Lagrange
% interpolants through them, which is what the solver actually evaluates.
% Markers are the nodes. Data from make_input_curves.jl in this directory,
% which converts to a load factor (p_L / 2.0 p.u.) and to cents/kWh.
%
\definecolor{cLoad}{HTML}{31748F}   % pine
\definecolor{cPrice}{HTML}{1F7A4D}  % cash green ($ price)
\definecolor{gridMajor}{HTML}{E8C4D4} % soft pink     -- major grid
\definecolor{gridMinor}{HTML}{D4D4E8} % soft lavender -- minor grid
%
\begin{figure}[!t]
\centering
\begin{tikzpicture}
\pgfplotsset{
  every axis/.append style={
    width=0.64\columnwidth, height=0.50\columnwidth, scale only axis,
    xmin=1, xmax=6,
    tick align=outside,
    tick label style={font=\footnotesize},
    label style={font=\footnotesize},
    axis line style={black, line width=0.6pt},
  },
}
% ---- Left axis: load factor, grids, full box ----
\begin{axis}[
    name=axL,
    ymin=0, ymax=1.08,
    xtick={1,2,3,4,5,6}, minor x tick num=1,
    ytick={0,0.2,0.4,0.6,0.8,1.0}, minor y tick num=1,
    xlabel={Time interval $t$},
    ylabel={Load factor $\lambda$},
    ylabel near ticks,
    grid=both,
    major grid style={line width=0.6pt, color=gridMajor, opacity=0.7},
    minor grid style={line width=0.4pt, color=gridMinor, opacity=0.6},
    legend style={
        at={(0.5,1.04)}, anchor=south,
        legend columns=2, draw=none, fill=none,
        column sep=1.0ex, font=\footnotesize,
    },
    legend image post style={line width=1.1pt},
]
% Load factor -- pine, solid, filled circles at the nodes
\addplot[cLoad, line width=1.5pt]
  table[col sep=comma, x=tau, y=lambda] {figures/input_curves.csv};
\addlegendentry{Load}
\addplot[only marks, mark=*, mark size=1.5pt, forget plot,
         mark options={fill=cLoad, draw=cLoad}]
  table[col sep=comma, x=tau, y=lambda] {figures/input_nodes.csv};
% Proxy legend entry for the price curve (plotted on the right axis below)
\addlegendimage{cPrice, dash dot, line width=1.5pt, mark=triangle*,
                mark size=1.7pt, mark options={fill=cPrice, draw=cPrice}}
\addlegendentry{Energy price}
\end{axis}
% ---- Right axis: energy price (cents/kWh), overlaid, right spine only ----
\begin{axis}[
    ymin=0, ymax=75,
    xmin=1, xmax=6,
    xtick=\empty,
    ytick={0,15,30,45,60,75},
    scaled y ticks=false,
    yticklabel style={/pgf/number format/fixed, /pgf/number format/precision=0},
    ylabel={Energy price [cents/kWh]},
    ylabel near ticks,
    axis x line=none,
    axis y line*=right,
    grid=none,
]
% Energy price -- cash green, dash-dot, filled triangles at the nodes
\addplot[cPrice, dash dot, line width=1.5pt]
  table[col sep=comma, x=tau, y=cents] {figures/input_curves.csv};
\addplot[only marks, mark=triangle*, mark size=1.7pt,
         mark options={fill=cPrice, draw=cPrice}]
  table[col sep=comma, x=tau, y=cents] {figures/input_nodes.csv};
\end{axis}
\end{tikzpicture}
\caption{Input profiles over the $T = 6$ horizon: load factor on the left axis,
energy price on the right, each specified at the $T$ horizon points.  Markers
are the data; the curves are the Lagrange interpolants~\eqref{eq:ocp_interp}
through them, and it is the curves, not the markers, that the solver evaluates.
The two agree exactly at every node, by construction, but between nodes they
need not be well behaved and already are not --- the load interpolant reaches
$0.998$ against a node maximum of $0.9$, and the price interpolant
$69.7$ cents/kWh against $62$.  Nothing samples those excursions here, because
the time-index state advances by exactly one per stage and so only ever lands on
a node; but the interpolant's \emph{derivatives} at the nodes are what the
backward pass consumes, and its degree is $T-1$.  This is the device that does
not survive a realistic horizon.}
\label{fig:input_curves}
\end{figure}
